\section{Intersections, cyclic commensurators, and rank bounds}
\label{pre:section}

We need three preliminary tools. The first controls intersections of
quasiconvex subgroups. The second produces an embedded Baumslag--Solitar
subgroup from a tree action. The third bounds the total reduced rank of
a family of conjugate intersections.

\subsection{Quasiconvex intersections}

All word metrics come from finite generating sets. A group is
\emph{word-hyperbolic} if geodesic triangles in a Cayley graph are
uniformly thin: each side lies in a fixed neighborhood of the other two.
A subgroup of a
hyperbolic group is \emph{quasiconvex} if every geodesic between its elements
lies in a fixed neighborhood of the subgroup. We use the Morse lemma and
its usual consequences: quasiconvex subgroups are finitely generated and
quasi-isometrically embedded, these two properties characterize
quasiconvexity, and infinite cyclic subgroups are quasiconvex. In particular,
quasiconvexity is transitive between hyperbolic groups, since
quasi-isometric embeddings compose. For these background facts, see
\cite[Theorem~III.H.1.7, Corollary~III.$\Gamma$.3.6(2), and
Corollary~III.$\Gamma$.3.10(1)]{BridsonHaefliger1999}.
For a subset \(Q\) of a hyperbolic
group, write \(\Lambda Q\) for its accumulation set in the boundary.

\begin{lemma}[Quasiconvex intersections]
\label{pre:quasiconvex-intersection}
Let \(K,L\) be quasiconvex subgroups of a finitely generated hyperbolic
group \(H\). Then \(K\cap L\) is quasiconvex and
\[
  \Lambda(K\cap L)=\Lambda K\cap\Lambda L.
\]
The same assertions hold for any finite family of quasiconvex subgroups.
\end{lemma}

\begin{proof}
First observe that, for fixed subgroups \(K,L\) of any finitely generated
group and \(R\geq 0\), there is a constant \(D\) such that
\begin{equation}
  N_R(K)\cap N_R(L)\subseteq N_D(K\cap L).
  \label{pre:neighborhood-intersection}
\end{equation}
Indeed, if \(x\) belongs to the left side, choose \(k\in K\) and \(l\in L\)
within distance \(R\) of \(x\). The element \(d=k^{-1}l\) lies in the finite
ball of radius \(2R\). For each value \(d\) that occurs, fix one pair
\(k_d\in K,l_d\in L\) with \(k_d^{-1}l_d=d\). Then
\[
  z=kk_d^{-1}=ll_d^{-1}\in K\cap L,
  \qquad d(x,z)\leq R+|k_d|.
\]
Taking a maximum over the finitely many \(d\) proves
\eqref{pre:neighborhood-intersection}.

A geodesic with endpoints in \(K\cap L\) lies in fixed neighborhoods of
both \(K\) and \(L\), so \eqref{pre:neighborhood-intersection} proves
quasiconvexity of their intersection.

Now let \(\xi\in\Lambda K\cap\Lambda L\). Quasiconvexity and limits of
geodesic segments show that a ray from the identity to \(\xi\) lies in
bounded neighborhoods of both subgroups. Applying
\eqref{pre:neighborhood-intersection} along the ray gives a sequence in
\(K\cap L\) converging to \(\xi\). Thus
\(\xi\in\Lambda(K\cap L)\). The reverse inclusion is immediate, and
induction gives the finite-family assertions.
These facts also appear in
\cite[Lemmas~2.6 and~2.7]{GitikMitraRipsSageev1998}.
\end{proof}

\begin{lemma}[Finite double-coset lists]
\label{pre:double-coset-finiteness}
Let \(K,L\) be quasiconvex subgroups of a finitely generated hyperbolic
group \(H\). There are only finitely many double cosets \(KgL\) such that
\(K\cap gLg^{-1}\) contains an infinite-order element. If \(H\) is
torsion-free, there are therefore only finitely many such double cosets
with nontrivial intersection.
\end{lemma}

\begin{proof}
Suppose that \(c\in K\cap gLg^{-1}\) has infinite order. Its two fixed
boundary points belong to both \(\Lambda K\) and \(\Lambda(gL)\): use the
sequences \(c^n\in K\) and \(c^ng\in gL\), for positive and negative \(n\).
A geodesic line between those points lies in a uniformly bounded
neighborhood of each of \(K\) and \(gL\). To see the uniformity, a
\(q\)-quasiconvex set contains finite geodesics between its points within
distance \(q\); passing to limits, and then comparing geodesics with the
same ideal endpoints, increases this bound by a constant depending only
on the hyperbolicity constant. The quasiconvexity constant of \(gL\)
equals that of \(L\).

At any point of this line choose nearby vertices \(k\in K\) and
\(gl\in gL\). Their distance is bounded by a constant depending only on
the two quasiconvexity constants and the hyperbolicity constant. Hence
\(k^{-1}gl\in KgL\) lies in a fixed finite ball. Every double coset under
consideration meets that ball, proving finiteness.
\end{proof}

The qualification about infinite-order elements is necessary in groups
with torsion. For example, the central subgroup \(C_2\) of the hyperbolic
group \(F_2\times C_2\) has nontrivial intersection with every conjugate,
but there are infinitely many \(C_2\)-double cosets.

For later applications to cosets, note also that
\[
  \Lambda(gL)=\Lambda(gLg^{-1}).
\]
Indeed, right multiplication by the fixed element \(g^{-1}\) moves each
point of \(gL\) a bounded distance.

\subsection{Commensuration in a tree action}

Two subgroups are \emph{commensurable} if their intersection has finite
index in both. For a subgroup \(C\leq H\), its commensurator is
\[
  \operatorname{Comm}_H(C)
  =\{h\in H: C\text{ and }hCh^{-1}\text{ are commensurable}\}.
\]
This is a subgroup. We first deduce the cyclic case directly from the
double-coset lemma.

\begin{lemma}
\label{pre:cyclic-commensurator}
If \(C\) is an infinite cyclic subgroup of a hyperbolic group \(H\), then
\(C\) has finite index in \(\operatorname{Comm}_H(C)\).
If \(H\) is torsion-free, this commensurator is infinite cyclic.
\end{lemma}

\begin{proof}
The subgroup \(C\) is quasiconvex. By
Lemma~\ref{pre:double-coset-finiteness}, its commensurator meets only
finitely many double cosets \(CgC\). For a representative \(g\) in the
commensurator, \(CgC\) is a union of
\([C:C\cap g^{-1}Cg]\) left \(C\)-cosets, a finite number. Thus the
commensurator is virtually cyclic. A torsion-free virtually cyclic group
is cyclic, giving the last assertion.
\end{proof}

All tree actions below are actions on simplicial trees without
inversions. An element is \emph{elliptic} if it fixes a vertex and
\emph{tree-hyperbolic} if it translates a bi-infinite geodesic.
We will need an embedded Baumslag--Solitar group, rather than just a
power-conjugacy relation. The following normal-form argument provides it.

\begin{lemma}[Graphs of cyclic groups]
\label{pre:gbs-subgroup}
Let \(Q\) be a finitely generated noncyclic group acting on a simplicial
tree without inversions. If every vertex and edge stabilizer is infinite
cyclic, then \(Q\) contains a subgroup isomorphic to \(BS(m,n)\) for some
nonzero integers \(m,n\).
\end{lemma}

\begin{proof}
Choose a vertex \(v\) and a finite symmetric generating set \(S\) of \(Q\).
The union of all \(Q\)-translates of the finite subtree
\[
  D=\bigcup_{s\in S}[v,sv]
\]
is connected: its segments connect the vertices in the orbit \(Qv\)
according to the Cayley graph of \(Q\). This union is exactly the convex
hull of \(Qv\), because its segments join orbit points. Restricting the
action to this hull gives a finite quotient, since every edge orbit has
a representative in \(D\). The retained vertex and edge stabilizers are
unchanged. Bass--Serre theory therefore realizes \(Q\) as the fundamental
group of a finite connected graph of infinite cyclic groups.

Collapse any non-loop edge whose map to one endpoint group is onto.
Such a collapse preserves the fundamental group and leaves a graph of
infinite cyclic groups. After finitely many collapses, every remaining
non-loop edge has both inclusion indices greater than one.

If a loop remains, the subgraph consisting of that loop and its vertex
has fundamental group
\[
  \langle x,t\mid t x^m t^{-1}=x^n\rangle=BS(m,n),
  \qquad m,n\neq 0.
\]
Its fundamental group injects into \(Q\) by the graph-of-groups normal
form theorem.

If a non-loop edge remains, its two-vertex subgraph likewise injects and
has fundamental group
\[
  \langle x,y\mid x^m=y^n\rangle,
  \qquad |m|,|n|>1.
\]
The element \(z=x^m=y^n\) is central and has infinite order. Modulo
\(\langle z\rangle\), this group is \(C_{|m|}*C_{|n|}\), in which the image
of \(xy\) has infinite order by free-product normal form. Thus the
commuting elements \(z\) and \(xy\) generate an embedded
\(\mathbb Z^2=BS(1,1)\).

Finally, if no edge remains, the graph has one vertex and \(Q\) is
cyclic, contrary to the hypothesis. This exhausts the possibilities.
\end{proof}

\begin{lemma}[Elliptic commensuration]
\label{pre:elliptic-commensurator}
Suppose that \(G\) acts on a simplicial tree \(\mathcal T\) without
inversions, with torsion-free hyperbolic vertex stabilizers. If an
infinite-order elliptic element \(a\in G\) has its cyclic subgroup
commensurated by a tree-hyperbolic element \(g\in G\), then \(G\) contains
a subgroup isomorphic to \(BS(m,n)\) for some nonzero integers \(m,n\).
\end{lemma}

\begin{proof}
Set \(L=\langle a,g\rangle\) and choose a vertex \(v\) fixed by \(a\).
Every element of \(L\) commensurates \(\langle a\rangle\). For each
\(h\in L\), commensurability supplies nonzero integers \(r,s\) with
\(a^r=ha^sh^{-1}\); hence a nonzero power of \(a\) fixes \(hv\).
A common multiple of finitely many such powers fixes any prescribed
finite subset of \(Lv\), and therefore fixes its convex hull pointwise.
It follows that every finite subtree of
\(\mathcal T_L=\operatorname{Hull}(Lv)\) is fixed by a nonzero power of
\(a\).

Let \(x\) be a vertex of \(\mathcal T_L\), and choose \(N>0\) with
\(a^N\in L_x\). Every element of \(L_x\) commensurates
\(\langle a^N\rangle\), now a subgroup of the vertex stabilizer \(G_x\).
Thus
\[
  \langle a^N\rangle
  \leq L_x
  \leq \operatorname{Comm}_{G_x}(\langle a^N\rangle).
\]
Lemma~\ref{pre:cyclic-commensurator} makes the last group infinite cyclic,
so \(L_x\) is infinite cyclic as well. Each edge stabilizer of
\(\mathcal T_L\) contains a nonzero power of \(a\) and is a subgroup of a
vertex stabilizer in \(L\); it too is infinite cyclic.

The group \(L\) is noncyclic. Otherwise, writing \(a=c^r\) and \(g=c^s\)
in a cyclic group, translation lengths would give
\(\ell(a)=|r|\ell(c)=0\) and hence \(\ell(g)=0\), contrary to
tree-hyperbolicity of \(g\). Apply Lemma~\ref{pre:gbs-subgroup} to the
action of the finitely generated group \(L\) on \(\mathcal T_L\).
\end{proof}

\subsection{Strong inertia in free groups}

We now formulate the rank estimate that will control noncyclic
intersections. For a finitely generated free group \(K\), put
\[
  \bar r(K)=\max\{\operatorname{rk}(K)-1,0\}.
\]
Strong inertia bounds the sum of reduced ranks of all conjugate
intersections; see
\cite{LintonHierarchies2025}.

\begin{definition}
\label{pre:strong-inertia-definition}
A finitely generated subgroup \(E\) of a finitely generated free group
\(P\) is \emph{strongly inert} if, for every finitely generated subgroup
\(S\leq P\),
\begin{equation}
  \sum_{SpE\in S\backslash P/E}
  \bar r(E\cap p^{-1}Sp)\leq\bar r(S).
  \label{pre:strong-inertia-inequality}
\end{equation}
The summands are independent of the chosen double-coset representatives.
Only finitely many intersections are nontrivial, by
Lemma~\ref{pre:double-coset-finiteness}.
\end{definition}

\begin{corollary}
\label{pre:cyclic-off-diagonal}
Let \(P\) be a finitely generated free group and \(E\leq P\) a strongly
inert subgroup. If \(g\in P\setminus E\), then
\(E\cap gEg^{-1}\) is trivial or cyclic.
\end{corollary}

\begin{proof}
Apply \eqref{pre:strong-inertia-inequality} with \(S=E\).
The identity double coset contributes \(\bar r(E)\), so all other
nonnegative summands vanish.
\end{proof}
